%-*-latex-*-
%-*-latex-*-
\newcommand\COURSE{ciss245}
\newcommand\ASSESSMENT{q2501}
\newcommand\ASSESSMENTTYPE{Quiz}
\newcommand\POINTS{\textwhite{xxx/xxx}}

\input{myquizpreamble}
\input{yliow}
\input{\COURSE}
%-*-latex-*-

%-*-latex-*-

%-*-latex-*-
\renewcommand\TITLE{\ASSESSMENTTYPE \ \ASSESSMENT}

\renewcommand\EMAIL{}
\renewcommand\AUTHOR{eagraham1@cougars.ccis.edu}

\textwidth=6in
\begin{document}
\topmatterthree

%------------------------------------------------------------------------------
\nextq
You want to write a random library that allows you to do this:
\begin{console}[fontsize=\footnotesize]
int i = randrange(5);           // i is a random integer in {0, 1, 2, 3, 4}
int j = randrange(3, 9);        // j is a random integer in {3, 4, 5, 6, 7, 8}
double x = randrange(5.0);      // x is a random double in interval [0.0, 5.0]
double y = randrange(2.0, 5.0); // y is a random double in interval [2.0, 5.0]
\end{console}
Write the header file below.
\\
\ANSWER
\begin{answercode}
// file: random.h
#ifndef RANDOM_H
#define RANDOM_H

int randrange(int end);
int randrange(int start, int end);
double randrange(double end);
double randrange(double start, double end);
void init_random();

#endif // RANDOM_H
\end{answercode}

%------------------------------------------------------------------------------
\nextq
Continuing the above question, implement the functions in above header file.
\\
\ANSWER
\begin{answercode}
// file: random.cpp
void init_random()
{
    std::srand((unsigned int) time(NULL));
}

int randrange(int end)
{
    return std::rand() % end;
}

int randrange(int start, int end) {
    return start + (std::rand() % (end - start));
}

double randrange(double end) {
    return (double(rand()) % end);
}

double randrange(double start, double end) {
    return start + double(std::rand()) % (end - start);
}
\end{answercode}

%------------------------------------------------------------------------------
%-*-latex-*-
\newpage
\input{instructions}
\end{document}

